\documentclass[11pt,a4paper]{article}
\usepackage[utf8]{inputenc}
\usepackage[T1]{fontenc}
\usepackage[english]{babel}
\usepackage[margin=2.3cm]{geometry}
\usepackage{xcolor}
\usepackage{booktabs}
\usepackage{longtable}
\usepackage{array}
\usepackage{amsmath,amssymb}
\usepackage{listings}
\usepackage[hidelinks]{hyperref}
\newcommand{\code}[1]{\texttt{#1}}
% Generated by tools/gen_concepts_values.py -- do not edit.
% Values come from mal-core/config/model_constants.yaml via the
% validated loader; a key that no longer exists makes \mskey fail
% the build with an error, in any document that references it.

\makeatletter
\newcommand{\mskey}[1]{%
  \expandafter\ifx\csname ms@k@#1\endcsname\relax
    \errmessage{Concept document: unknown model constant key `#1'}%
  \else
    \csname ms@k@#1\endcsname%
  \fi}
\makeatother

\expandafter\def\csname ms@k@adults.daily_mort_basal\endcsname{0.87}
\expandafter\def\csname ms@k@adults.mort_cap\endcsname{0.93}
\expandafter\def\csname ms@k@adults.mort_floor\endcsname{0.6}
\expandafter\def\csname ms@k@adults.opt_c\endcsname{25.0}
\expandafter\def\csname ms@k@adults.sigma_c\endcsname{15.0}
\expandafter\def\csname ms@k@aquatic_stages.egg_daily_mort\endcsname{0.08}
\expandafter\def\csname ms@k@aquatic_stages.larva_bh_alpha\endcsname{0.05}
\expandafter\def\csname ms@k@aquatic_stages.larva_bh_s0\endcsname{0.95}
\expandafter\def\csname ms@k@aquatic_stages.last_larval_instar\endcsname{4}
\expandafter\def\csname ms@k@aquatic_stages.pupa_daily_mort\endcsname{0.05}
\expandafter\def\csname ms@k@aquatic_stages.salinity_missing_psu\endcsname{0.0}
\expandafter\def\csname ms@k@aquatic_stress.desicc_base_daily\endcsname{0.1}
\expandafter\def\csname ms@k@aquatic_stress.desicc_grace_days\endcsname{5.0}
\expandafter\def\csname ms@k@aquatic_stress.desicc_ramp_decay\endcsname{0.3}
\expandafter\def\csname ms@k@aquatic_stress.vulnerable_instar_max\endcsname{2}
\expandafter\def\csname ms@k@aquatic_stress.wash_fraction_max\endcsname{0.6}
\expandafter\def\csname ms@k@aquatic_stress.wash_threshold_mm\endcsname{40.0}
\expandafter\def\csname ms@k@capacity.persistent_availability\endcsname{1.0}
\expandafter\def\csname ms@k@contracts.cogs_state\endcsname{"4.0"}
\expandafter\def\csname ms@k@contracts.env_nc\endcsname{"3.0"}
\expandafter\def\csname ms@k@contracts.habitat\endcsname{"4.0"}
\expandafter\def\csname ms@k@dispersal.flight_hour\endcsname{20}
\expandafter\def\csname ms@k@dispersal.night_end_hour\endcsname{6}
\expandafter\def\csname ms@k@dispersal.night_start_hour\endcsname{18}
\expandafter\def\csname ms@k@dispersal.snap_radius_m\endcsname{500.0}
\expandafter\def\csname ms@k@dispersal.wind_flight_hours\endcsname{4.0}
\expandafter\def\csname ms@k@dispersal.wind_migration_prob\endcsname{0.05}
\expandafter\def\csname ms@k@dispersal.wind_self_speed_ms\endcsname{1.0}
\expandafter\def\csname ms@k@dispersal.wind_survival\endcsname{0.85}
\expandafter\def\csname ms@k@eip.base_c\endcsname{16.0}
\expandafter\def\csname ms@k@eip.threshold_gd\endcsname{110.0}
\expandafter\def\csname ms@k@geodesy.earth_radius_m\endcsname{6371008.8}
\expandafter\def\csname ms@k@geodesy.metres_per_deg_lat\endcsname{111195.08023353292}
\expandafter\def\csname ms@k@gonotrophic.cycle_duration_days\endcsname{2.0}
\expandafter\def\csname ms@k@gonotrophic.egg_batch_max\endcsname{170}
\expandafter\def\csname ms@k@gonotrophic.egg_batch_mean\endcsname{120}
\expandafter\def\csname ms@k@gonotrophic.egg_batch_min\endcsname{30}
\expandafter\def\csname ms@k@gonotrophic.egg_batch_n\endcsname{240}
\expandafter\def\csname ms@k@gonotrophic.feeding_success_rate\endcsname{0.825}
\expandafter\def\csname ms@k@gonotrophic.first_cycle_days\endcsname{4.0}
\expandafter\def\csname ms@k@gonotrophic.resting_duration_days\endcsname{0.5}
\expandafter\def\csname ms@k@habitat.activation_min_engine\endcsname{0.01}
\expandafter\def\csname ms@k@habitat.candidate_min_engine\endcsname{0.01}
\expandafter\def\csname ms@k@habitat_types.dynamics.crop_pluvial\endcsname{"pluvial"}
\expandafter\def\csname ms@k@habitat_types.dynamics.open_pluvial\endcsname{"pluvial"}
\expandafter\def\csname ms@k@habitat_types.dynamics.surface_persistent\endcsname{"persistent"}
\expandafter\def\csname ms@k@habitat_types.dynamics.urban_persistent\endcsname{"persistent"}
\expandafter\def\csname ms@k@habitat_types.dynamics.wetland_persistent\endcsname{"persistent"}
\expandafter\def\csname ms@k@habitat_types.dynamics.wetland_pluvial\endcsname{"pluvial"}
\expandafter\def\csname ms@k@habitat_types.types\endcsname{"surface\_persistent, wetland\_persistent, urban\_persistent, wetland\_pluvial, open\_pluvial, crop\_pluvial"}
\expandafter\def\csname ms@k@host_seeking.indoor_attraction_add\endcsname{0.72}
\expandafter\def\csname ms@k@host_seeking.indoor_fraction_rural\endcsname{0.3}
\expandafter\def\csname ms@k@host_seeking.indoor_fraction_urban\endcsname{0.72}
\expandafter\def\csname ms@k@host_seeking.starvation_days\endcsname{7}
\expandafter\def\csname ms@k@host_seeking.urbanicity_add\endcsname{0.2}
\expandafter\def\csname ms@k@ingest.candidate_min_ingest\endcsname{1e-06}
\expandafter\def\csname ms@k@ingest.cell_size_hint_m\endcsname{1000.0}
\expandafter\def\csname ms@k@ingest.coastline_buffer_m\endcsname{2000.0}
\expandafter\def\csname ms@k@ingest.coastline_resolution\endcsname{"f"}
\expandafter\def\csname ms@k@ingest.fill_frac_valid_max\endcsname{0.5}
\expandafter\def\csname ms@k@ingest.grid_degrees_step_max\endcsname{1}
\expandafter\def\csname ms@k@ingest.land_mask_threshold\endcsname{0.5}
\expandafter\def\csname ms@k@ingest.min_dilation_radius_cells\endcsname{1}
\expandafter\def\csname ms@k@ingest.nc_compression_level\endcsname{4}
\expandafter\def\csname ms@k@ingest.nodata_sentinel\endcsname{-9999.0}
\expandafter\def\csname ms@k@ingest.nodata_sentinel_int\endcsname{-9999}
\expandafter\def\csname ms@k@ingest.ocean_mask_threshold\endcsname{0.5}
\expandafter\def\csname ms@k@ingest.residential_from_indoor_frac\endcsname{0.8}
\expandafter\def\csname ms@k@ingest.sea_water_floor\endcsname{0.05}
\expandafter\def\csname ms@k@ingest.valid_count_epsilon\endcsname{1e-09}
\expandafter\def\csname ms@k@integration.directed_flight_budget_h\endcsname{0.4856}
\expandafter\def\csname ms@k@integration.dt_olfactory_s\endcsname{0.25}
\expandafter\def\csname ms@k@integration.dt_terminal_s\endcsname{0.04}
\expandafter\def\csname ms@k@integration.landing_speed_hold_s\endcsname{3.0}
\expandafter\def\csname ms@k@integration.landing_speed_mm_s\endcsname{50.0}
\expandafter\def\csname ms@k@integration.phase_duration_hours\endcsname{"12.0, 4.0, 6.0, 2.0"}
\expandafter\def\csname ms@k@integration.plume_memory_s\endcsname{0.7}
\expandafter\def\csname ms@k@mating.female_attendance_rate\endcsname{0.1033}
\expandafter\def\csname ms@k@mating.local_swarm_male_offer_factor\endcsname{0.0}
\expandafter\def\csname ms@k@mating.male_females_per_night_max\endcsname{3}
\expandafter\def\csname ms@k@mating.male_move_capacity_m\endcsname{50.0}
\expandafter\def\csname ms@k@mating.male_swarm_reference_size\endcsname{100.0}
\expandafter\def\csname ms@k@mating.p_mate_male_fraction\endcsname{0.072}
\expandafter\def\csname ms@k@mating.p_mate_pairs_per_individual\endcsname{0.167}
\expandafter\def\csname ms@k@mating.p_mate_size_exponent\endcsname{0.0}
\expandafter\def\csname ms@k@mating.rain_suppresses_swarming\endcsname{true}
\expandafter\def\csname ms@k@mating.swarm_duration_min\endcsname{28.3}
\expandafter\def\csname ms@k@mating.swarm_height_m\endcsname{2.87}
\expandafter\def\csname ms@k@mating.swarm_markers_per_km2\endcsname{1.0}
\expandafter\def\csname ms@k@mating.swarm_markers_source\endcsname{"engineering prior"}
\expandafter\def\csname ms@k@mating.swarm_start_min_after_sunset\endcsname{5.0}
\expandafter\def\csname ms@k@oviposition.detect_range_m\endcsname{2000.0}
\expandafter\def\csname ms@k@oviposition.min_attraction\endcsname{0.3}
\expandafter\def\csname ms@k@oviposition.scale_m\endcsname{50.0}
\expandafter\def\csname ms@k@oviposition.step_m\endcsname{50.0}
\expandafter\def\csname ms@k@oviposition.timeout_days\endcsname{7}
\expandafter\def\csname ms@k@pluvial.cold_start_z\endcsname{0.0}
\expandafter\def\csname ms@k@pluvial.delta_days\endcsname{1.0}
\expandafter\def\csname ms@k@pluvial.dry_memory_days\endcsname{5.0}
\expandafter\def\csname ms@k@pluvial.rain_scale_mm\endcsname{2.5}
\expandafter\def\csname ms@k@rng.aquatic_seed_larva_mortality\endcsname{123}
\expandafter\def\csname ms@k@rng.aquatic_seed_stage_mortality\endcsname{42}
\expandafter\def\csname ms@k@rng.aquatic_seed_washout\endcsname{789}
\expandafter\def\csname ms@k@species.arabiensis.activity_weights\endcsname{"0.02, 0.1, 0.8, 0.08"}
\expandafter\def\csname ms@k@species.arabiensis.aestivation_supported\endcsname{true}
\expandafter\def\csname ms@k@species.arabiensis.anemotaxis_deviation_deg\endcsname{30.0}
\expandafter\def\csname ms@k@species.arabiensis.anemotaxis_rule\endcsname{"upwind"}
\expandafter\def\csname ms@k@species.arabiensis.arrested_daily_survival\endcsname{0.986}
\expandafter\def\csname ms@k@species.arabiensis.blood_meal_refractory_days\endcsname{0.5}
\expandafter\def\csname ms@k@species.arabiensis.blood_meal_size_ul\endcsname{0.3}
\expandafter\def\csname ms@k@species.arabiensis.blood_meal_target_ul\endcsname{0.5}
\expandafter\def\csname ms@k@species.arabiensis.casting_angle_deg\endcsname{90.0}
\expandafter\def\csname ms@k@species.arabiensis.casting_probability\endcsname{0.1}
\expandafter\def\csname ms@k@species.arabiensis.co2_activation_radius_m\endcsname{100.0}
\expandafter\def\csname ms@k@species.arabiensis.daily_move_capacity_m\endcsname{590.0}
\expandafter\def\csname ms@k@species.arabiensis.directed_flight_budget_h\endcsname{0.4856}
\expandafter\def\csname ms@k@species.arabiensis.eggs_per_ul\endcsname{223.0}
\expandafter\def\csname ms@k@species.arabiensis.exploration_range_m\endcsname{100.0}
\expandafter\def\csname ms@k@species.arabiensis.exploration_step_correlation\endcsname{0.5}
\expandafter\def\csname ms@k@species.arabiensis.exploration_step_m\endcsname{50.0}
\expandafter\def\csname ms@k@species.arabiensis.first_cycle_target_ul\endcsname{0.8}
\expandafter\def\csname ms@k@species.arabiensis.flight_speed_ms\endcsname{0.3375}
\expandafter\def\csname ms@k@species.arabiensis.gonotrophic_arrest_prob_dry\endcsname{0.4454}
\expandafter\def\csname ms@k@species.arabiensis.male_daily_survival\endcsname{0.773}
\expandafter\def\csname ms@k@species.arabiensis.male_ratio\endcsname{0.5}
\expandafter\def\csname ms@k@species.arabiensis.meal_size_min_ul\endcsname{0.1}
\expandafter\def\csname ms@k@species.arabiensis.olfactory_radius_m\endcsname{50.0}
\expandafter\def\csname ms@k@species.arabiensis.plume_memory_s\endcsname{0.7}
\expandafter\def\csname ms@k@species.arabiensis.pref.cattle\endcsname{0.2}
\expandafter\def\csname ms@k@species.arabiensis.pref.goat\endcsname{0.15}
\expandafter\def\csname ms@k@species.arabiensis.pref.human\endcsname{0.5}
\expandafter\def\csname ms@k@species.arabiensis.pref.sheep\endcsname{0.1}
\expandafter\def\csname ms@k@species.arabiensis.pref.wildlife\endcsname{0.05}
\expandafter\def\csname ms@k@species.arabiensis.probing_success_rate\endcsname{0.7275}
\expandafter\def\csname ms@k@species.arabiensis.salinity_hi_tol_ppt\endcsname{15.0}
\expandafter\def\csname ms@k@species.arabiensis.salinity_opt_ppt\endcsname{0.0}
\expandafter\def\csname ms@k@species.arabiensis.salinity_slope\endcsname{0.15}
\expandafter\def\csname ms@k@species.arabiensis.terminal_radius_m\endcsname{0.15}
\expandafter\def\csname ms@k@species.arabiensis.visual_radius_m\endcsname{0.485}
\expandafter\def\csname ms@k@species.arabiensis.weights.crop_pluvial\endcsname{0.1}
\expandafter\def\csname ms@k@species.arabiensis.weights.open_pluvial\endcsname{0.55}
\expandafter\def\csname ms@k@species.arabiensis.weights.surface_persistent\endcsname{0.1}
\expandafter\def\csname ms@k@species.arabiensis.weights.urban_persistent\endcsname{0.15}
\expandafter\def\csname ms@k@species.arabiensis.weights.wetland_persistent\endcsname{0.1}
\expandafter\def\csname ms@k@species.arabiensis.weights.wetland_pluvial\endcsname{0.0}
\expandafter\def\csname ms@k@species.arabiensis.wind_block_threshold_ms\endcsname{0.3375}
\expandafter\def\csname ms@k@species.arabiensis.wind_reference\endcsname{"u100/v100"}
\expandafter\def\csname ms@k@species.coluzzii.activity_weights\endcsname{"0.0, 0.255, 0.48, 0.265"}
\expandafter\def\csname ms@k@species.coluzzii.aestivation_supported\endcsname{true}
\expandafter\def\csname ms@k@species.coluzzii.anemotaxis_deviation_deg\endcsname{30.0}
\expandafter\def\csname ms@k@species.coluzzii.anemotaxis_rule\endcsname{"upwind"}
\expandafter\def\csname ms@k@species.coluzzii.arrested_daily_survival\endcsname{0.986}
\expandafter\def\csname ms@k@species.coluzzii.blood_meal_refractory_days\endcsname{0.5}
\expandafter\def\csname ms@k@species.coluzzii.blood_meal_size_ul\endcsname{0.3}
\expandafter\def\csname ms@k@species.coluzzii.blood_meal_target_ul\endcsname{0.5}
\expandafter\def\csname ms@k@species.coluzzii.casting_angle_deg\endcsname{90.0}
\expandafter\def\csname ms@k@species.coluzzii.casting_probability\endcsname{0.1}
\expandafter\def\csname ms@k@species.coluzzii.co2_activation_radius_m\endcsname{100.0}
\expandafter\def\csname ms@k@species.coluzzii.daily_move_capacity_m\endcsname{590.0}
\expandafter\def\csname ms@k@species.coluzzii.directed_flight_budget_h\endcsname{0.4856}
\expandafter\def\csname ms@k@species.coluzzii.eggs_per_ul\endcsname{223.0}
\expandafter\def\csname ms@k@species.coluzzii.exploration_range_m\endcsname{100.0}
\expandafter\def\csname ms@k@species.coluzzii.exploration_step_correlation\endcsname{0.5}
\expandafter\def\csname ms@k@species.coluzzii.exploration_step_m\endcsname{50.0}
\expandafter\def\csname ms@k@species.coluzzii.first_cycle_target_ul\endcsname{0.8}
\expandafter\def\csname ms@k@species.coluzzii.flight_speed_ms\endcsname{0.3375}
\expandafter\def\csname ms@k@species.coluzzii.gonotrophic_arrest_prob_dry\endcsname{0.2704}
\expandafter\def\csname ms@k@species.coluzzii.male_daily_survival\endcsname{0.773}
\expandafter\def\csname ms@k@species.coluzzii.male_ratio\endcsname{0.5}
\expandafter\def\csname ms@k@species.coluzzii.meal_size_min_ul\endcsname{0.1}
\expandafter\def\csname ms@k@species.coluzzii.olfactory_radius_m\endcsname{50.0}
\expandafter\def\csname ms@k@species.coluzzii.plume_memory_s\endcsname{0.7}
\expandafter\def\csname ms@k@species.coluzzii.pref.cattle\endcsname{0.005}
\expandafter\def\csname ms@k@species.coluzzii.pref.goat\endcsname{0.003}
\expandafter\def\csname ms@k@species.coluzzii.pref.human\endcsname{0.99}
\expandafter\def\csname ms@k@species.coluzzii.pref.sheep\endcsname{0.001}
\expandafter\def\csname ms@k@species.coluzzii.pref.wildlife\endcsname{0.001}
\expandafter\def\csname ms@k@species.coluzzii.probing_success_rate\endcsname{0.7275}
\expandafter\def\csname ms@k@species.coluzzii.salinity_hi_tol_ppt\endcsname{30.0}
\expandafter\def\csname ms@k@species.coluzzii.salinity_opt_ppt\endcsname{0.0}
\expandafter\def\csname ms@k@species.coluzzii.salinity_slope\endcsname{0.08}
\expandafter\def\csname ms@k@species.coluzzii.terminal_radius_m\endcsname{0.15}
\expandafter\def\csname ms@k@species.coluzzii.visual_radius_m\endcsname{0.485}
\expandafter\def\csname ms@k@species.coluzzii.weights.crop_pluvial\endcsname{0.05}
\expandafter\def\csname ms@k@species.coluzzii.weights.open_pluvial\endcsname{0.25}
\expandafter\def\csname ms@k@species.coluzzii.weights.surface_persistent\endcsname{0.4}
\expandafter\def\csname ms@k@species.coluzzii.weights.urban_persistent\endcsname{0.2}
\expandafter\def\csname ms@k@species.coluzzii.weights.wetland_persistent\endcsname{0.1}
\expandafter\def\csname ms@k@species.coluzzii.weights.wetland_pluvial\endcsname{0.0}
\expandafter\def\csname ms@k@species.coluzzii.wind_block_threshold_ms\endcsname{0.3375}
\expandafter\def\csname ms@k@species.coluzzii.wind_reference\endcsname{"u100/v100"}
\expandafter\def\csname ms@k@species.default_species\endcsname{"coluzzii"}
\expandafter\def\csname ms@k@species.funestus.activity_weights\endcsname{"0.02, 0.1, 0.8, 0.08"}
\expandafter\def\csname ms@k@species.funestus.aestivation_supported\endcsname{false}
\expandafter\def\csname ms@k@species.funestus.anemotaxis_deviation_deg\endcsname{30.0}
\expandafter\def\csname ms@k@species.funestus.anemotaxis_rule\endcsname{"upwind"}
\expandafter\def\csname ms@k@species.funestus.arrested_daily_survival\endcsname{0.986}
\expandafter\def\csname ms@k@species.funestus.blood_meal_refractory_days\endcsname{0.5}
\expandafter\def\csname ms@k@species.funestus.blood_meal_size_ul\endcsname{0.3}
\expandafter\def\csname ms@k@species.funestus.blood_meal_target_ul\endcsname{0.5}
\expandafter\def\csname ms@k@species.funestus.casting_angle_deg\endcsname{90.0}
\expandafter\def\csname ms@k@species.funestus.casting_probability\endcsname{0.1}
\expandafter\def\csname ms@k@species.funestus.co2_activation_radius_m\endcsname{100.0}
\expandafter\def\csname ms@k@species.funestus.daily_move_capacity_m\endcsname{590.0}
\expandafter\def\csname ms@k@species.funestus.directed_flight_budget_h\endcsname{0.4856}
\expandafter\def\csname ms@k@species.funestus.eggs_per_ul\endcsname{223.0}
\expandafter\def\csname ms@k@species.funestus.exploration_range_m\endcsname{100.0}
\expandafter\def\csname ms@k@species.funestus.exploration_step_correlation\endcsname{0.5}
\expandafter\def\csname ms@k@species.funestus.exploration_step_m\endcsname{50.0}
\expandafter\def\csname ms@k@species.funestus.first_cycle_target_ul\endcsname{0.8}
\expandafter\def\csname ms@k@species.funestus.flight_speed_ms\endcsname{0.3375}
\expandafter\def\csname ms@k@species.funestus.gonotrophic_arrest_prob_dry\endcsname{0.0}
\expandafter\def\csname ms@k@species.funestus.male_daily_survival\endcsname{0.773}
\expandafter\def\csname ms@k@species.funestus.male_ratio\endcsname{0.5}
\expandafter\def\csname ms@k@species.funestus.meal_size_min_ul\endcsname{0.1}
\expandafter\def\csname ms@k@species.funestus.olfactory_radius_m\endcsname{50.0}
\expandafter\def\csname ms@k@species.funestus.plume_memory_s\endcsname{0.7}
\expandafter\def\csname ms@k@species.funestus.pref.cattle\endcsname{0.005}
\expandafter\def\csname ms@k@species.funestus.pref.goat\endcsname{0.003}
\expandafter\def\csname ms@k@species.funestus.pref.human\endcsname{0.99}
\expandafter\def\csname ms@k@species.funestus.pref.sheep\endcsname{0.001}
\expandafter\def\csname ms@k@species.funestus.pref.wildlife\endcsname{0.001}
\expandafter\def\csname ms@k@species.funestus.probing_success_rate\endcsname{0.7275}
\expandafter\def\csname ms@k@species.funestus.salinity_hi_tol_ppt\endcsname{15.0}
\expandafter\def\csname ms@k@species.funestus.salinity_opt_ppt\endcsname{0.0}
\expandafter\def\csname ms@k@species.funestus.salinity_slope\endcsname{0.15}
\expandafter\def\csname ms@k@species.funestus.terminal_radius_m\endcsname{0.15}
\expandafter\def\csname ms@k@species.funestus.visual_radius_m\endcsname{0.485}
\expandafter\def\csname ms@k@species.funestus.weights.crop_pluvial\endcsname{0.05}
\expandafter\def\csname ms@k@species.funestus.weights.open_pluvial\endcsname{0.1}
\expandafter\def\csname ms@k@species.funestus.weights.surface_persistent\endcsname{0.45}
\expandafter\def\csname ms@k@species.funestus.weights.urban_persistent\endcsname{0.1}
\expandafter\def\csname ms@k@species.funestus.weights.wetland_persistent\endcsname{0.3}
\expandafter\def\csname ms@k@species.funestus.weights.wetland_pluvial\endcsname{0.0}
\expandafter\def\csname ms@k@species.funestus.wind_block_threshold_ms\endcsname{0.3375}
\expandafter\def\csname ms@k@species.funestus.wind_reference\endcsname{"u100/v100"}
\expandafter\def\csname ms@k@species.gambiae_ss.activity_weights\endcsname{"0.02, 0.1, 0.8, 0.08"}
\expandafter\def\csname ms@k@species.gambiae_ss.aestivation_supported\endcsname{true}
\expandafter\def\csname ms@k@species.gambiae_ss.anemotaxis_deviation_deg\endcsname{30.0}
\expandafter\def\csname ms@k@species.gambiae_ss.anemotaxis_rule\endcsname{"upwind"}
\expandafter\def\csname ms@k@species.gambiae_ss.arrested_daily_survival\endcsname{0.986}
\expandafter\def\csname ms@k@species.gambiae_ss.blood_meal_refractory_days\endcsname{0.5}
\expandafter\def\csname ms@k@species.gambiae_ss.blood_meal_size_ul\endcsname{0.3}
\expandafter\def\csname ms@k@species.gambiae_ss.blood_meal_target_ul\endcsname{0.5}
\expandafter\def\csname ms@k@species.gambiae_ss.casting_angle_deg\endcsname{90.0}
\expandafter\def\csname ms@k@species.gambiae_ss.casting_probability\endcsname{0.1}
\expandafter\def\csname ms@k@species.gambiae_ss.co2_activation_radius_m\endcsname{100.0}
\expandafter\def\csname ms@k@species.gambiae_ss.daily_move_capacity_m\endcsname{590.0}
\expandafter\def\csname ms@k@species.gambiae_ss.directed_flight_budget_h\endcsname{0.4856}
\expandafter\def\csname ms@k@species.gambiae_ss.eggs_per_ul\endcsname{223.0}
\expandafter\def\csname ms@k@species.gambiae_ss.exploration_range_m\endcsname{100.0}
\expandafter\def\csname ms@k@species.gambiae_ss.exploration_step_correlation\endcsname{0.5}
\expandafter\def\csname ms@k@species.gambiae_ss.exploration_step_m\endcsname{50.0}
\expandafter\def\csname ms@k@species.gambiae_ss.first_cycle_target_ul\endcsname{0.8}
\expandafter\def\csname ms@k@species.gambiae_ss.flight_speed_ms\endcsname{0.3375}
\expandafter\def\csname ms@k@species.gambiae_ss.gonotrophic_arrest_prob_dry\endcsname{0.6838}
\expandafter\def\csname ms@k@species.gambiae_ss.male_daily_survival\endcsname{0.773}
\expandafter\def\csname ms@k@species.gambiae_ss.male_ratio\endcsname{0.5}
\expandafter\def\csname ms@k@species.gambiae_ss.meal_size_min_ul\endcsname{0.1}
\expandafter\def\csname ms@k@species.gambiae_ss.olfactory_radius_m\endcsname{50.0}
\expandafter\def\csname ms@k@species.gambiae_ss.plume_memory_s\endcsname{0.7}
\expandafter\def\csname ms@k@species.gambiae_ss.pref.cattle\endcsname{0.005}
\expandafter\def\csname ms@k@species.gambiae_ss.pref.goat\endcsname{0.003}
\expandafter\def\csname ms@k@species.gambiae_ss.pref.human\endcsname{0.99}
\expandafter\def\csname ms@k@species.gambiae_ss.pref.sheep\endcsname{0.001}
\expandafter\def\csname ms@k@species.gambiae_ss.pref.wildlife\endcsname{0.001}
\expandafter\def\csname ms@k@species.gambiae_ss.probing_success_rate\endcsname{0.7275}
\expandafter\def\csname ms@k@species.gambiae_ss.salinity_hi_tol_ppt\endcsname{20.0}
\expandafter\def\csname ms@k@species.gambiae_ss.salinity_opt_ppt\endcsname{0.0}
\expandafter\def\csname ms@k@species.gambiae_ss.salinity_slope\endcsname{0.12}
\expandafter\def\csname ms@k@species.gambiae_ss.terminal_radius_m\endcsname{0.15}
\expandafter\def\csname ms@k@species.gambiae_ss.visual_radius_m\endcsname{0.485}
\expandafter\def\csname ms@k@species.gambiae_ss.weights.crop_pluvial\endcsname{0.1}
\expandafter\def\csname ms@k@species.gambiae_ss.weights.open_pluvial\endcsname{0.65}
\expandafter\def\csname ms@k@species.gambiae_ss.weights.surface_persistent\endcsname{0.05}
\expandafter\def\csname ms@k@species.gambiae_ss.weights.urban_persistent\endcsname{0.15}
\expandafter\def\csname ms@k@species.gambiae_ss.weights.wetland_persistent\endcsname{0.05}
\expandafter\def\csname ms@k@species.gambiae_ss.weights.wetland_pluvial\endcsname{0.0}
\expandafter\def\csname ms@k@species.gambiae_ss.wind_block_threshold_ms\endcsname{0.3375}
\expandafter\def\csname ms@k@species.gambiae_ss.wind_reference\endcsname{"u100/v100"}
\expandafter\def\csname ms@k@species.k_ref\endcsname{50.0}
\expandafter\def\csname ms@k@species.melas.activity_weights\endcsname{"0.02, 0.1, 0.8, 0.08"}
\expandafter\def\csname ms@k@species.melas.aestivation_supported\endcsname{false}
\expandafter\def\csname ms@k@species.melas.anemotaxis_deviation_deg\endcsname{30.0}
\expandafter\def\csname ms@k@species.melas.anemotaxis_rule\endcsname{"upwind"}
\expandafter\def\csname ms@k@species.melas.arrested_daily_survival\endcsname{0.986}
\expandafter\def\csname ms@k@species.melas.blood_meal_refractory_days\endcsname{0.5}
\expandafter\def\csname ms@k@species.melas.blood_meal_size_ul\endcsname{0.3}
\expandafter\def\csname ms@k@species.melas.blood_meal_target_ul\endcsname{0.5}
\expandafter\def\csname ms@k@species.melas.casting_angle_deg\endcsname{90.0}
\expandafter\def\csname ms@k@species.melas.casting_probability\endcsname{0.1}
\expandafter\def\csname ms@k@species.melas.co2_activation_radius_m\endcsname{100.0}
\expandafter\def\csname ms@k@species.melas.daily_move_capacity_m\endcsname{590.0}
\expandafter\def\csname ms@k@species.melas.directed_flight_budget_h\endcsname{0.4856}
\expandafter\def\csname ms@k@species.melas.eggs_per_ul\endcsname{223.0}
\expandafter\def\csname ms@k@species.melas.exploration_range_m\endcsname{100.0}
\expandafter\def\csname ms@k@species.melas.exploration_step_correlation\endcsname{0.5}
\expandafter\def\csname ms@k@species.melas.exploration_step_m\endcsname{50.0}
\expandafter\def\csname ms@k@species.melas.first_cycle_target_ul\endcsname{0.8}
\expandafter\def\csname ms@k@species.melas.flight_speed_ms\endcsname{0.3375}
\expandafter\def\csname ms@k@species.melas.gonotrophic_arrest_prob_dry\endcsname{0.0}
\expandafter\def\csname ms@k@species.melas.male_daily_survival\endcsname{0.773}
\expandafter\def\csname ms@k@species.melas.male_ratio\endcsname{0.5}
\expandafter\def\csname ms@k@species.melas.meal_size_min_ul\endcsname{0.1}
\expandafter\def\csname ms@k@species.melas.olfactory_radius_m\endcsname{50.0}
\expandafter\def\csname ms@k@species.melas.plume_memory_s\endcsname{0.7}
\expandafter\def\csname ms@k@species.melas.pref.cattle\endcsname{0.005}
\expandafter\def\csname ms@k@species.melas.pref.goat\endcsname{0.003}
\expandafter\def\csname ms@k@species.melas.pref.human\endcsname{0.99}
\expandafter\def\csname ms@k@species.melas.pref.sheep\endcsname{0.001}
\expandafter\def\csname ms@k@species.melas.pref.wildlife\endcsname{0.001}
\expandafter\def\csname ms@k@species.melas.probing_success_rate\endcsname{0.7275}
\expandafter\def\csname ms@k@species.melas.salinity_hi_tol_ppt\endcsname{38.0}
\expandafter\def\csname ms@k@species.melas.salinity_opt_ppt\endcsname{0.0}
\expandafter\def\csname ms@k@species.melas.salinity_slope\endcsname{0.02}
\expandafter\def\csname ms@k@species.melas.terminal_radius_m\endcsname{0.15}
\expandafter\def\csname ms@k@species.melas.visual_radius_m\endcsname{0.485}
\expandafter\def\csname ms@k@species.melas.weights.crop_pluvial\endcsname{0.05}
\expandafter\def\csname ms@k@species.melas.weights.open_pluvial\endcsname{0.4}
\expandafter\def\csname ms@k@species.melas.weights.surface_persistent\endcsname{0.25}
\expandafter\def\csname ms@k@species.melas.weights.urban_persistent\endcsname{0.15}
\expandafter\def\csname ms@k@species.melas.weights.wetland_persistent\endcsname{0.15}
\expandafter\def\csname ms@k@species.melas.weights.wetland_pluvial\endcsname{0.0}
\expandafter\def\csname ms@k@species.melas.wind_block_threshold_ms\endcsname{0.3375}
\expandafter\def\csname ms@k@species.melas.wind_reference\endcsname{"u100/v100"}
\expandafter\def\csname ms@k@species.neutral.weights.crop_pluvial\endcsname{0.05}
\expandafter\def\csname ms@k@species.neutral.weights.open_pluvial\endcsname{0.4}
\expandafter\def\csname ms@k@species.neutral.weights.surface_persistent\endcsname{0.25}
\expandafter\def\csname ms@k@species.neutral.weights.urban_persistent\endcsname{0.15}
\expandafter\def\csname ms@k@species.neutral.weights.wetland_persistent\endcsname{0.15}
\expandafter\def\csname ms@k@species.neutral.weights.wetland_pluvial\endcsname{0.0}
\expandafter\def\csname ms@k@species.stephensi.activity_weights\endcsname{"0.02, 0.1, 0.8, 0.08"}
\expandafter\def\csname ms@k@species.stephensi.aestivation_supported\endcsname{false}
\expandafter\def\csname ms@k@species.stephensi.anemotaxis_deviation_deg\endcsname{30.0}
\expandafter\def\csname ms@k@species.stephensi.anemotaxis_rule\endcsname{"upwind"}
\expandafter\def\csname ms@k@species.stephensi.arrested_daily_survival\endcsname{0.986}
\expandafter\def\csname ms@k@species.stephensi.blood_meal_refractory_days\endcsname{0.5}
\expandafter\def\csname ms@k@species.stephensi.blood_meal_size_ul\endcsname{0.3}
\expandafter\def\csname ms@k@species.stephensi.blood_meal_target_ul\endcsname{0.5}
\expandafter\def\csname ms@k@species.stephensi.casting_angle_deg\endcsname{90.0}
\expandafter\def\csname ms@k@species.stephensi.casting_probability\endcsname{0.1}
\expandafter\def\csname ms@k@species.stephensi.co2_activation_radius_m\endcsname{100.0}
\expandafter\def\csname ms@k@species.stephensi.daily_move_capacity_m\endcsname{590.0}
\expandafter\def\csname ms@k@species.stephensi.directed_flight_budget_h\endcsname{0.4856}
\expandafter\def\csname ms@k@species.stephensi.eggs_per_ul\endcsname{223.0}
\expandafter\def\csname ms@k@species.stephensi.exploration_range_m\endcsname{100.0}
\expandafter\def\csname ms@k@species.stephensi.exploration_step_correlation\endcsname{0.5}
\expandafter\def\csname ms@k@species.stephensi.exploration_step_m\endcsname{50.0}
\expandafter\def\csname ms@k@species.stephensi.first_cycle_target_ul\endcsname{0.8}
\expandafter\def\csname ms@k@species.stephensi.flight_speed_ms\endcsname{0.3375}
\expandafter\def\csname ms@k@species.stephensi.gonotrophic_arrest_prob_dry\endcsname{0.0}
\expandafter\def\csname ms@k@species.stephensi.male_daily_survival\endcsname{0.773}
\expandafter\def\csname ms@k@species.stephensi.male_ratio\endcsname{0.5}
\expandafter\def\csname ms@k@species.stephensi.meal_size_min_ul\endcsname{0.1}
\expandafter\def\csname ms@k@species.stephensi.olfactory_radius_m\endcsname{50.0}
\expandafter\def\csname ms@k@species.stephensi.plume_memory_s\endcsname{0.7}
\expandafter\def\csname ms@k@species.stephensi.pref.cattle\endcsname{0.15}
\expandafter\def\csname ms@k@species.stephensi.pref.goat\endcsname{0.05}
\expandafter\def\csname ms@k@species.stephensi.pref.human\endcsname{0.7}
\expandafter\def\csname ms@k@species.stephensi.pref.sheep\endcsname{0.05}
\expandafter\def\csname ms@k@species.stephensi.pref.wildlife\endcsname{0.05}
\expandafter\def\csname ms@k@species.stephensi.probing_success_rate\endcsname{0.7275}
\expandafter\def\csname ms@k@species.stephensi.salinity_hi_tol_ppt\endcsname{20.0}
\expandafter\def\csname ms@k@species.stephensi.salinity_opt_ppt\endcsname{0.0}
\expandafter\def\csname ms@k@species.stephensi.salinity_slope\endcsname{0.15}
\expandafter\def\csname ms@k@species.stephensi.terminal_radius_m\endcsname{0.15}
\expandafter\def\csname ms@k@species.stephensi.visual_radius_m\endcsname{0.485}
\expandafter\def\csname ms@k@species.stephensi.weights.crop_pluvial\endcsname{0.05}
\expandafter\def\csname ms@k@species.stephensi.weights.open_pluvial\endcsname{0.4}
\expandafter\def\csname ms@k@species.stephensi.weights.surface_persistent\endcsname{0.25}
\expandafter\def\csname ms@k@species.stephensi.weights.urban_persistent\endcsname{0.15}
\expandafter\def\csname ms@k@species.stephensi.weights.wetland_persistent\endcsname{0.15}
\expandafter\def\csname ms@k@species.stephensi.weights.wetland_pluvial\endcsname{0.0}
\expandafter\def\csname ms@k@species.stephensi.wind_block_threshold_ms\endcsname{0.3375}
\expandafter\def\csname ms@k@species.stephensi.wind_reference\endcsname{"u100/v100"}
\expandafter\def\csname ms@k@thermal.efd_max_eggs_per_day\endcsname{0.5}
\expandafter\def\csname ms@k@thermal.efd_tmax_c\endcsname{34.0}
\expandafter\def\csname ms@k@thermal.efd_tmin_c\endcsname{16.0}
\expandafter\def\csname ms@k@thermal.efd_topt_c\endcsname{25.0}
\expandafter\def\csname ms@k@thermal.egg_development_scale\endcsname{0.00062}
\expandafter\def\csname ms@k@thermal.egg_development_tmax_c\endcsname{40.0}
\expandafter\def\csname ms@k@thermal.egg_development_tmin_c\endcsname{10.0}
\expandafter\def\csname ms@k@thermal.larva_development_scale\endcsname{0.001}
\expandafter\def\csname ms@k@thermal.larva_development_tmax_c\endcsname{40.0}
\expandafter\def\csname ms@k@thermal.larva_development_tmin_c\endcsname{10.0}
\expandafter\def\csname ms@k@thermal.pupa_development_scale\endcsname{0.00066}
\expandafter\def\csname ms@k@thermal.pupa_development_tmax_c\endcsname{40.0}
\expandafter\def\csname ms@k@thermal.pupa_development_tmin_c\endcsname{12.0}
\expandafter\def\csname ms@k@transmission.beta_hv\endcsname{0.4}
\expandafter\def\csname ms@k@transmission.beta_vh\endcsname{0.5}
\expandafter\def\csname ms@k@transmission.human_immunity_days\endcsname{180}
\expandafter\def\csname ms@k@transmission.human_incubation_days\endcsname{12}
\expandafter\def\csname ms@k@transmission.human_infectious_days\endcsname{20}
\expandafter\def\csname ms@k@transmission.initial_human_prevalence\endcsname{0.05}
\expandafter\def\csname ms@k@transmission.initial_vector_infected_frac\endcsname{0.0}
\expandafter\def\csname ms@k@urban.beta_urban\endcsname{0.75}
\expandafter\def\csname ms@k@urban.rain_pools_built_damage\endcsname{0.3}
\expandafter\def\csname ms@k@version\endcsname{3}
\expandafter\def\csname ms@k@wildlife.class_suitability.bare\endcsname{0.2}
\expandafter\def\csname ms@k@wildlife.class_suitability.built\endcsname{0.05}
\expandafter\def\csname ms@k@wildlife.class_suitability.crops\endcsname{0.3}
\expandafter\def\csname ms@k@wildlife.class_suitability.flooded_vegetation\endcsname{0.45}
\expandafter\def\csname ms@k@wildlife.class_suitability.grass\endcsname{0.6}
\expandafter\def\csname ms@k@wildlife.class_suitability.shrub_and_scrub\endcsname{0.7}
\expandafter\def\csname ms@k@wildlife.class_suitability.snow_and_ice\endcsname{0.1}
\expandafter\def\csname ms@k@wildlife.class_suitability.trees\endcsname{0.9}
\expandafter\def\csname ms@k@wildlife.class_suitability.water\endcsname{0.0}
\expandafter\def\csname ms@k@wildlife.weights.habitat\endcsname{0.5}
\expandafter\def\csname ms@k@wildlife.weights.remoteness\endcsname{0.2}
\expandafter\def\csname ms@k@wildlife.weights.water\endcsname{0.3}


\title{The \emph{scoring} stage of the MalariaSentinel pipeline\\[3pt]
\large A conceptual description of how a finished simulation is turned into a scorecard}
\author{MalariaSentinel conceptual-documentation layer}
\date{}

\begin{document}
\maketitle

\begin{abstract}
The \emph{scoring} stage runs after an agent-based-model execution has finished and
answers one question: \emph{is this simulation behaving like the malaria
entomology it is meant to emulate, and did my last change make it better or
worse?} It does so by reading the artifacts of a \emph{finished} run --- the
state COGs, the cohort and aquatic logs, the transmission log --- and by
comparing what it finds against bands taken from the entomological
literature. Its output is a single JSON file, the \emph{scorecard}, written
inside the run directory. The scorecard carries a weighted geometric mean
over the composite dimensions, a set of binary engineering gates reported
apart from that mean, and an explicit list of what could not be evaluated and
why. This document describes what the stage measures, how it computes what it
computes, what it deliberately does not do, and where the current code and the
current description of the system disagree.
\end{abstract}

\section{Purpose, and what the stage is not for}

\subsection{The question the stage answers}

The stage exists to answer a comparison question about a \emph{finished} run.
The registry module states the intent in its own opening paragraph: the
weights it carries ``reflect validation value for the question
\emph{is the system good?}'', and the weights are set so that external
validation and core biology dominate while engineering invariants are demoted
to gates rather than counted as dimensions
(\code{mal-core/src/mal\_core/scoring/registry.py:1-18}). Concretely, the
stage answers two coupled questions:

\begin{enumerate}
\item \emph{Is this run in the right biological neighbourhood?} Each
      dimension asks one narrow, checkable question --- do adult lifespans
      fall in the field-observed window, does the aquatic population have
      the stage structure implied by development durations, does the adult
      population spread as far as mark--recapture studies say it should ---
      and returns a number between zero and one.
\item \emph{Did my change help?} Because every dimension is measured the same
      way on every run, two scorecards are comparable, and the feedback step
      can report the change in the composite and the dimensions that moved
      backwards.
\end{enumerate}

``Having answered it'' means: a scorecard exists, inside the run directory,
with one entry per registered scorer stating what was measured, what the
target was, whether the dimension passed, and --- for everything that could
not be measured --- why not
(\code{mal-core/src/mal\_core/scoring/runner.py:120-146}). A scorecard that
reports nothing but skips is a valid, and honest, answer: it says the run
could not be evaluated, and which artifacts were missing.

\subsection{The negative purpose}

The stage is \emph{not}:

\begin{itemize}
\item \textbf{A re-run.} The stage neither executes nor re-parameterises the
      model. Its only contract with the simulation is one method per scorer,
      which receives a run directory and a context dictionary and returns a
      result
      (\code{mal-core/src/mal\_core/scoring/scorers/base.py:78-79}).
      Producing the run is the ABM stage's job.
\item \textbf{A decision gate.} The \code{score} command computes the
      scorecard, writes it and prints a summary; it does not fail the run
      and does not block anything downstream
      (\code{mal-core/src/mal\_core/cli.py:409-463}). Nothing in the working
      tree reads the composite and refuses to proceed. The verdict label is
      a human-readable reading of the number, not an enforced criterion.
      The one non-zero exit in the command is an argument error and has
      nothing to do with the verdict: a missing \code{--run-dir} prints a
      message and raises \code{typer.Exit(code=1)}
      (\code{mal-core/src/mal\_core/cli.py:448-450}); once the run directory
      has been supplied, the command always writes the scorecard and exits
      zero (\code{cli.py:455-463}).
\item \textbf{A place where the language model decides.} The LLM judge's
      verdict is \emph{advisory}: it can only add an optional key to the
      scorecard, and it can never change the composite or the verdict
      (\code{mal-core/src/mal\_core/scoring/runner.py:135-142}).
\item \textbf{The old calibration harness.} There is no pytest calibration
      harness in the tree: no \code{calibration/} directory, no
      \code{thresholds.yaml}, no diff-versus-best tracker. The harness and
      its duplicated scorers were deleted, and the base module says so
      (\code{mal-core/src/mal\_core/scoring/scorers/base.py:1-6}). The
      thresholds that used to live there are now inlined in each scorer
      module, or read from the shared model configuration where they belong
      to the model.
\item \textbf{A producer of runs.} Two sibling commands exist that do
      produce a run before scoring --- \code{validate-detections} and
      \code{validate-cases} --- but they belong to the measurement
      campaigns described in \S\ref{sec:campaigns}, not to the stage itself.
\end{itemize}

\subsection{How to read this document}

The stance note (\code{docs/concepts/README.md}) asks for three things this
document tries to do consistently. First, every magnitude is named with its
identity and its units, and every number is either generated from the shared
model configuration or written here with its \emph{provenance tier} stated.
Second, choices are named as choices: the AOI slug, the weights, the radius,
the horizons --- each is something a reader can vary. Third, every claim
resolves to a file and a line a reviewer can open, and where the code does
not settle a question the document says so rather than guessing.

The provenance tiers used throughout are:

\begin{longtable}{@{}p{2.6cm}p{11.4cm}@{}}
\toprule
Tier & Meaning \\
\midrule
\endhead
\textbf{Model definition} & The value is part of the model's definition and
lives in the shared configuration file
(\code{mal-core/config/model\_constants.yaml}); it is written here through the
generated macro \code{\string\mskey} and can be changed by changing the
model. \\
\textbf{Literature-anchored} & The value is a band or a target taken from
published field or semi-field work; it lives as a literal in the scorer
module, with the source named in that module's docstring. \\
\textbf{Engineering prior} & The value is a choice made by the implementer
to make a check computable (sampling caps, floors, tie tolerances). It has
no literature anchor and is not derived from the model. \\
\textbf{Scenario default} & The value is a default that only holds for one
particular AOI, campaign or run configuration; it is stated in the scorer
that uses it. \\
\bottomrule
\end{longtable}

Values that fall into the first tier are never typed by hand in this
document: they are pulled in through the generated macros. Values in the
other three tiers are literals in the code, and the evidence appendix lists
them all with the file and line where each one is read.

The tier is a classification made by \emph{this document}, not a fact read
from the code. Every value that is not pulled from the shared configuration
is a plain literal in its scorer module, and the code draws no distinction
between a value that is a scenario default and one that is an engineering
prior --- both are bare literals of the same kind
(\code{d16\_detection\_coverage.py:43-45},
\code{d4\_stability.py:56-66}). Where a value could sit in either of those
two tiers, the tier given below is this document's judgement, and a reviewer
should read it as such.

% Generated by tools/gen_concepts_values.py -- do not edit.
\paragraph{The output contracts}
\begin{longtable}{@{}llllll@{}}
\toprule
Artefact & Producer & Consumers & Version & Bands / variables & Sentinel \\
\midrule
\endhead
\code{env} & \code{mal\_core.ingest.env.build\_env\_tensor} & abm, prediction, scoring & "3.0" & rainfall, air\_temp\_c, jrc\_persistent, lc\_built, lc\_wetland, lc\_open\_plus\_water, lc\_crop, land\_frac (optional), salinity\_ppt (optional) & none \\
\code{habitat} & \code{mal\_core.ingest.env.build\_env\_tensor} & abm &  & hab\_type, habitat\_potential, row, col & none \\
\code{state\_cog} & \code{mal\_abm\_fast::write\_state\_cog} & scoring, prediction, training & "4.0" & adult\_patch\_count, adult\_cell\_count & -9999.0 \\
\code{state\_sidecar} & \code{mal\_abm\_fast::write\_state\_sidecar} & scoring & "4.0" &  &  \\
\code{manifest} & \code{mal\_core.download.manifest.update\_dataset} & download, ingest, abm &  &  &  \\
\code{host\_static} & \code{mal\_core.ingest.hosts.build\_host\_dataset} & abm &  & humans, livestock, urban, buildings &  \\
\code{mobility\_day} & \code{mal\_core.ingest.mobility.build\_mobility\_dataset} & abm &  &  &  \\
\bottomrule
\end{longtable}

\paragraph{The state sidecar's keys, in emission order}
\begin{longtable}{@{}l@{}}
\toprule
Key \\
\midrule
\endhead
\code{crs} \\
\code{transform} \\
\code{aoi\_slug} \\
\code{scale} \\
\code{year} \\
\code{month} \\
\code{seed} \\
\code{n\_rollouts} \\
\code{rollout\_index} \\
\code{generator\_version} \\
\code{abm\_params\_hash} \\
\code{contract\_version} \\
\code{band\_names} \\
\code{nodata} \\
\code{shape} \\
\bottomrule
\end{longtable}


\section{The scorecard: the stage's one output}

\subsection{What goes in}

The stage's input is \emph{one directory}: the output directory of a
finished ABM run
(\code{mal-core/src/mal\_core/scoring/runner.py:3-7}). The runner refuses a
path that is not a directory
(\code{mal-core/src/mal\_core/scoring/runner.py:66-67}). Alongside the
directory the caller may supply an AOI slug, a scoring configuration file,
and three name filters
(\code{mal-core/src/mal\_core/scoring/runner.py:54-63}).

The directory is not read by a fixed manifest; it is \emph{discovered} by
globbing, with a documented fallback chain per artifact class
(\code{mal-core/src/mal\_core/scoring/scorers/base.py:82-161}):

\begin{longtable}{@{}p{3.4cm}p{8.6cm}p{3.2cm}@{}}
\toprule
Artifact & How it is found & Primary consumers \\
\midrule
\endhead
State COGs (one per snapshot) & \code{*\_day*.tif} excluding any name
containing \code{transmission}; failing that \code{state\_day*.tif}, then
\code{state\_seed*.tif}, then a single \code{state.tif}
(\code{base.py:89-102}) & D1, D5, D8, D15, D16, D21, G6, G24 \\
Cohort log (adult counts per day) & \code{*\_cohort.json}, then
\code{cohort.json}, then \code{cohort\_seed*.json}
(\code{base.py:116-125}) & D2, D3, D4, D16 (extinction check) \\
Aquatic logs (larval stages per day) & \code{*\_day*\_aquatic.json}, falling
back to the summary \code{*\_aquatic.json}
(\code{base.py:128-139}) & D11, D15 \\
Run summary (contract version, AOI, grid) & the first \code{*.json} in the
directory that is not named like a cohort/aquatic/day/scorecard file and
that carries a \code{contract\_version} or \code{aoi\_slug} key
(\code{base.py:142-161}) & the context; D16, D21, D25 \\
Transmission log (daily incidence) & \code{*\_transmission\_daily.json}
(\code{d25\_cases\_ghana.py:135}) & D25 \\
\end{longtable}

Three artifacts are read but are \emph{not} in the run directory. They are the
AOI's own data, resolved by convention: the runner substitutes the AOI slug
into repository-relative glob patterns and takes the first existing hit
(\code{base.py:204-223}). The occurrence table
(\code{data/\{aoi\}/\{aoi\}\_occurrence.txt}) feeds D16
(\code{d16\_detection\_coverage.py:47-50}); the host grid
(\code{data/\{aoi\}/\{aoi\}\_host\_static.nc}) feeds G12
(\code{gates.py:161-164}); the urban class raster
(\code{data/\{aoi\}/\{aoi\}\_urban\_class.tif}) feeds G24
(\code{gates.py:255-265}). One scorer downloads its own data: D25 fetches a
pinned DHIMS-II CSV (URL and SHA-256 fixed in the module) into
\code{data/ghana/scoring/} when the file is absent or its checksum does not
match (\code{d25\_cases\_ghana.py:39-77}).

\subsection{What comes out}

The output is one file, \code{scorecard.json}, written \emph{inside} the run
directory (\code{mal-core/src/mal\_core/scoring/runner.py:120-146}). Its keys
are:

\begin{longtable}{@{}p{3.6cm}p{11.4cm}@{}}
\toprule
Key & Content \\
\midrule
\endhead
\code{run\_dir} & The scored directory, as given. \\
\code{aoi} & The AOI slug, or \code{null} when none was supplied. \\
\code{composite} & The weighted geometric mean over the contributing composite
dimensions, rounded to four decimals. \\
\code{verdict} & One of \code{VIABLE}, \code{BORDERLINE},
\code{REGRESSED}, \code{COLLAPSED}: a label on the composite. \\
\code{scores} & One entry per scorer that was run, holding its score, the
measured value, the target string, its diagnostics, its pass flag and its
status. \\
\code{composite\_dimensions} & The dimensions that contributed to the
composite, each with its score and its weight. \\
\code{gates} & The gates that ran, each with the same entry shape. \\
\code{skipped} & A map from scorer name to the reason it was not evaluated. \\
\code{timestamp} & The UTC ISO-8601 instant the card was written. \\
\code{verdict\_llm} & Present only when the language-model judge succeeded;
absent otherwise. \\
\bottomrule
\end{longtable}

\subsection{Who consumes it, and under what contract}

The scorecard is written for a human and for the feedback step. The
composite is rounded to four decimals, so a consumer that needs to compare
two runs should compare the stored values rather than recompute from the
raw diagnostics. There are three consumers in the tree:

\begin{itemize}
\item The command-line summary, which prints the composite and verdict, one
      line per composite dimension (marked \code{LOW} when the dimension did
      not reach its own threshold), the gates, and the skips
      (\code{runner.py:149-177}).
\item The feedback renderer, which turns a scorecard --- optionally against a
      baseline scorecard --- into a Markdown report
      (\code{mal-core/src/mal\_core/scoring/feedback.py:16-53}).
\item The language-model judge, which is handed the finished scorecard
      exactly as it will be written to disk
      (\code{mal-core/src/mal\_core/scoring/runner.py:135-142}).
\end{itemize}

No other stage of the pipeline reads the scorecard. The downstream stages
consume the run's own artifacts --- the state COG and the sidecar, under
contract \mskey{contracts.cogs_state} --- not the scorecard. That is worth
saying plainly: \emph{the scorecard is a report about a run, not an input to
the next stage}.

\section{The registry: how a scorer is declared}

\subsection{The scorer contract}

Every scorer in the stage is a subclass of the abstract base class
\code{Scorer}
(\code{mal-core/src/mal\_core/scoring/scorers/base.py:56-79}). A scorer
declares a name, a kind, a description, a minimum score, a hard floor and a
manual flag, and implements one method,
\code{score(run\_dir, ctx) -> ScorerResult}
(\code{mal-core/src/mal\_core/scoring/scorers/base.py:78-79}):

\begin{longtable}{@{}p{2.9cm}p{11.1cm}@{}}
\toprule
Member & Meaning \\
\midrule
\endhead
\code{name} & The registry key, e.g. \code{D2\_survival}: the name used on the
command line, in the configuration file and in the scorecard. \\
\code{kind} & Either \code{"composite"} or \code{"gate"}. The kind decides
whether the result enters the composite or is reported apart from it. \\
\code{description} & One line of prose, shown by \code{malariasim score
--list} (\code{mal-core/src/mal\_core/cli.py:440-446}). \\
\code{min\_score} & The threshold the scorer compares its own score against
to set its \emph{passed} flag. \\
\code{hard\_floor} & A second declared threshold. It is declared by eleven
dimensions and read by no code in the package (see
\S\ref{sec:absences}). \\
\code{manual} & Marks a scorer that must never run by default because it
needs per-AOI data (\code{base.py:72-76}). \\
\end{longtable}

The result type carries six fields (\code{base.py:26-43}): the \code{score}
that would enter the composite, the \code{value} measured in its own units, a
human-readable \code{target} band, free-form \code{diagnostics}, the
\code{passed} flag and the \code{status}, which is either \code{ok} or
\code{skipped}. A skipped result is built by a dedicated factory that zeroes
the score and the value, clears \code{passed} and records the reason in
\code{diagnostics["skipped\_reason"]} (\code{base.py:45-54}).

\subsection{The registry is code, not a data file}

\code{build\_registry()} instantiates every scorer and gate in the package
and returns a mapping from name to a frozen record holding the class and its
kind (\code{mal-core/src/mal\_core/scoring/registry.py:47-51},
\code{registry.py:79-104}). There is no registry file on disk, so a scorer
cannot be declared in one place and implemented in another, and the registry
cannot drift out of date with respect to the code.

Adding a scorer therefore takes three steps, and all three are needed:

\begin{enumerate}
\item write the module under \code{mal-core/src/mal\_core/scoring/scorers/}
      with a \code{Scorer} subclass; the package's export surface is
      \code{scorers/\_\_init\_\_.py:1-47};
\item instantiate it in \code{build\_registry()} (\code{registry.py:80-100});
\item if it is a composite dimension, give it a weight in
      \code{COMPOSITE\_WEIGHTS} (\code{registry.py:53-67}).
\end{enumerate}

A scorer that is registered but given no weight still runs, because the
configuration falls back to a weight of one
(\code{mal-core/src/mal\_core/scoring/config.py:30-31}); the listing, however,
reads the weight table directly and would show it with weight zero
(\code{registry.py:116}).

\subsection{Weights, and why they are what they are}

The weight table is the single source of truth for the composite:

% Generated by tools/gen_concepts_values.py -- do not edit.
\paragraph{Composite dimensions and their weights}
\begin{longtable}{@{}llrl@{}}
\toprule
Dimension & Weight & Min. score & Hard floor \\
\midrule
\endhead
\code{D11\_aquatic\_structure} & 1.0 & 0.7 & 0.3 \\
\code{D15\_persistence} & 3.0 & 0.7 & 0.4 \\
\code{D16\_detection\_coverage} & 4.0 & 0.5 & 0.2 \\
\code{D19\_pool\_persistence} & 1.0 & 0.7 & 0.0 \\
\code{D1\_habitat\_occupancy} & 1.0 & 0.5 & 0.2 \\
\code{D20\_washout\_response} & 1.0 & 0.7 & 0.0 \\
\code{D21\_dispersal\_kernel} & 1.0 & 0.7 & 0.3 \\
\code{D25\_cases\_ghana} & 2.0 & 0.5 & 0.2 \\
\code{D2\_survival} & 4.0 & 0.5 & 0.2 \\
\code{D3\_eip} & 3.0 & 0.3 & 0.1 \\
\code{D4\_stability} & 2.0 & 0.7 & 0.3 \\
\code{D5\_morans} & 1.0 & 0.5 & 0.2 \\
\code{D8\_coupling} & 2.0 & 0.4 & 0.1 \\
\bottomrule
\end{longtable}

\paragraph{Gates (evaluated separately from the composite)}
\begin{longtable}{@{}ll@{}}
\toprule
Gate & Description \\
\midrule
\endhead
\code{G6\_data\_valid} & All state bands finite and >= 0 (COG/NaN write integrity) \\
\code{G7\_determinism} & Same-seed state files hash-identical \\
\code{G10\_perf} & Wall time <= 30 s (when run\_metadata.json records it) \\
\code{G12\_host\_input} & Input human grid total \textasciitilde{} census reference (input validation) \\
\code{G14\_mobility} & Mobility OD matrix rows sum to 1.0 (input validation) \\
\code{G24\_urban\_ratio} & Urban/rural density ratio <= 0.7 (capacity-scaling guardrail) \\
\bottomrule
\end{longtable}


The registry's opening paragraph states the reasoning: the weights reflect
validation value for the question \emph{is the system good?}, so external
validation and core biology dominate, while engineering invariants are
demoted to gates rather than weighted into the biology
(\code{registry.py:1-6}). The two heaviest dimensions are the
detection-coverage dimension and the adult-survival dimension; the
one-weight family is the set of occupancy, spatial-autocorrelation,
aquatic-stage, habitat-persistence, washout and dispersal dimensions.

\subsection{Enabling, disabling, filtering}

By default every scorer that is \emph{not} manual is enabled
(\code{config.py:34-48}). A configuration file may enable or disable a
scorer and override a weight; its shape is a \code{scorers} mapping whose
entries carry \code{enabled} and \code{weight} keys
(\code{config.py:80-89}).

The command line is applied on top of the file in this order: \code{--enable}
first, then the configuration file, then \code{--only} (intersect with the
enabled set) and \code{--skip} (subtract from it)
(\code{config.py:51-95}). An unknown scorer name raises a
\code{ValueError} naming the offending string (\code{config.py:68-72}); the
command-line lists are comma-separated (\code{cli.py:452-454}).

The \emph{manual} flag is the stage's honesty mechanism for AOI-dependent
scorers: a manual scorer needs per-AOI data that not every AOI has --- an
occurrence table, an urban-class raster, a case series --- so it stays out
of the default set until an operator who knows the data is there activates
it explicitly (\code{base.py:72-76}). Three scorers carry the flag: the
detection-coverage dimension
(\code{d16\_detection\_coverage.py:60}), the reported-cases dimension
(\code{d25\_cases\_ghana.py:184}) and the urban-ratio gate
(\code{gates.py:259}).

\section{Order of operations, and what that order decides}

The runner executes a fixed sequence
(\code{mal-core/src/mal\_core/scoring/runner.py:54-146}):

\begin{enumerate}
\item \textbf{Validate the run directory.} A path that is not a directory
      raises \code{NotADirectoryError} before anything else happens
      (\code{runner.py:66-67}).
\item \textbf{Resolve the configuration.} The file, \code{--only},
      \code{--skip} and \code{--enable} are folded into one enabled set and
      one weight map (\code{runner.py:68-76}).
\item \textbf{Build the registry and the context.} The context carries the
      AOI slug, the run summary metadata and the run directory path
      (\code{runner.py:78-82}, \code{runner.py:34-40}).
\item \textbf{Execute the scorers, in alphabetical order of their names}
      (\code{runner.py:84}). Each call is wrapped in a \code{try} block and
      an exception becomes a skipped entry carrying the exception text
      (\code{runner.py:96-98}); the code states the reason --- a broken
      scorer must not sink the run.
\item \textbf{Triage the results into three buckets}: skipped results go to
      \code{skipped}, composite results to the composite map, gate results
      to \code{gates} (\code{runner.py:99-106}). A scorer that was enabled
      but not run because configuration disabled it is also recorded in
      \code{skipped}, with a different reason (\code{runner.py:86-93}).
\item \textbf{Compute the composite} over the composite map
      (\code{runner.py:108-118}).
\item \textbf{Label the verdict} from the composite (\code{runner.py:24-31}).
\item \textbf{Call the optional judge}, attaching its reply only if it
      succeeded (\code{runner.py:135-142}).
\item \textbf{Write the scorecard} into the run directory
      (\code{runner.py:144-146}).
\end{enumerate}

\subsection{Consequences worth stating}

\begin{itemize}
\item \textbf{A skipped dimension shrinks the composite's meaning; it does
      not punish it.} A dimension that cannot be measured removes itself
      from both the numerator and the weight sum (\code{runner.py:112-118}).
      Adding a dimension therefore narrows what a composite covers without
      invalidating older composites --- and two scorecards computed over
      different dimension sets are not directly comparable.
\item \textbf{A crashing scorer is invisible in the composite.} Because the
      exception is converted into a skip, a systematic failure in one
      dimension appears as a missing dimension rather than as a low score.
      The reason string in \code{skipped} is the only trace.
\item \textbf{Gates are computed but cannot affect the composite.} The two
      maps are disjoint by construction (\code{runner.py:103-106}).
\item \textbf{The judge cannot affect anything.} It is called after the
      composite and the verdict are final, and its reply is attached under a
      separate key (\code{runner.py:135-142}).
\item \textbf{Writing happens last.} A judge failure is absorbed inside the
      judge module, so the scorecard is written either way
      (\code{llm\_scorer.py:359-368}).
\item \textbf{Ordering is alphabetical, not dependency-aware.} Scorers do
      not depend on one another; their only shared inputs are the files on
      disk.
\end{itemize}

\section{The composite: one number from many}
\label{sec:composite}

\subsection{The formula}

The composite is a \emph{weighted geometric mean} over the composite
dimensions that produced an \emph{ok} result
(\code{mal-core/src/mal\_core/scoring/runner.py:108-118}):

\begin{equation}
C = \exp\!\left( \frac{ \sum_i w_i \ln s_i }{ \sum_i w_i } \right),
\qquad s_i \in (0,1],
\quad w_i = \texttt{weight\_of}(\text{name}_i)
\end{equation}

The weight of a dimension comes from the configuration, which falls back to
the registry table and then to one (\code{config.py:30-31}).

The geometric mean is the right shape for this job, and its properties are
the stage's failure semantics. Because $\ln s \to -\infty$ as $s \to 0$:

\begin{itemize}
\item \textbf{One zero collapses everything.} The composite is computed only
      when \emph{all} contributing scores are strictly positive; otherwise it
      is zero (\code{runner.py:111}). A run in which one dimension collapsed
      is a run whose composite is zero, which is the intended reading: a
      collapsed dimension is not something to be averaged away.
\item \textbf{No contributing dimension also gives zero}
      (\code{runner.py:110-111}). A consumer must therefore distinguish
      \emph{the model scored zero} from \emph{nothing could be scored} by
      reading \code{composite\_dimensions} and \code{skipped}.
\item \textbf{Low dimensions dominate high ones.} A single dimension far
      below the others pulls the composite well under the arithmetic mean,
      because the geometric mean is dominated by its smallest term. The
      composite is as much a statement about the worst dimension as about
      the average one.
\end{itemize}

\subsection{What a dimension's score means}

Each dimension returns a number in $[0,1]$ answering a question of the form
\emph{how far is this measurement from the expected band?}. Six shapes
occur among the thirteen, and they are not interchangeable --- a score of
one means something slightly different in each of them:

\begin{itemize}
\item \textbf{Plateau and falloff.} The score is one inside an accepted band
      and decays as a Gaussian away from it --- for example the occupancy
      dimension (\code{d1\_habitat\_occupancy.py:62-66}) and the incidence
      dimension (\code{d3\_eip.py:54-57}).
\item \textbf{Plateau and falloff, averaged over two bands.} The dimension
      computes one such quantity per band and averages the two ---
      \code{d4\_stability.py:56-68},
      \code{d11\_aquatic\_structure.py:90-105},
      \code{d21\_dispersal\_kernel.py:115-126}.
\item \textbf{Gaussian without a plateau.} The survival dimension has no
      band branch: its score is an unconditional Gaussian centred on a mean
      lifespan of $11$ days with a width of $4$ days
      (\code{d2\_survival.py:55}), so the score is one only at exactly
      $11$ days, and the ends of the stated $8$--$15$ day window do not
      score one.
\item \textbf{Ratio against a target.} The score is the measured quantity
      divided by its target, and is one when the target is met --- the
      coupling dimension divides the correlation by $0.50$ and caps at one
      (\code{d8\_coupling.py:54}), the detection-coverage dimension divides
      the coverage by its target (\code{d16\_detection\_coverage.py:128}).
\item \textbf{Geometric mean of three sub-scores.} The persistence
      dimension combines its survival, no-extinction and expansion
      sub-scores by a geometric mean
      (\code{d15\_persistence.py:123-126}).
\item \textbf{Weighted sum of boolean checks.} The score is the sum of the
      weights of the checks that hold --- the two contract dimensions
      (\code{d19\_pool\_persistence.py:126-134},
      \code{d20\_washout\_response.py:65-73}).
\end{itemize}

A dimension score is therefore neither a probability nor a percentage: it is
\emph{the degree to which one measurement of one run matches an expected
band}, and it is interpretable only together with the \code{value} and
\code{target} fields that accompany it.

\subsection{The minimum score, the hard floor, and the pass flag}

\code{min\_score} is a per-dimension declaration: the scorer compares its
own score against it and sets the \emph{passed} flag (\code{base.py:70}; for
example \code{d2\_survival.py:64}). The generated table above lists, for
each dimension, its weight, its minimum score and its hard floor.

The pattern is not universal, and two exceptions are worth naming. The
coupling dimension never compares its score against its declared minimum
score at all: it declares \code{min\_score = 0.4}
(\code{d8\_coupling.py:25}) but sets its flag from the raw correlation
instead, \code{passed = r >= 0.30} (\code{d8\_coupling.py:60}), so its
declared minimum score is inert. The two contract dimensions set their
flags from a literal $0.70$ rather than from \code{min\_score}
(\code{d19\_pool\_persistence.py:144},
\code{d20\_washout\_response.py:83}); there the literal coincides with the
declared minimum score, so the difference has no behavioural effect.

One consequence is worth stating explicitly: \emph{a dimension that fails
its minimum score still enters the composite}. The runner puts a dimension
into the composite map whenever its status is \emph{ok}, regardless of
whether it passed (\code{runner.py:100-106}). The pass flag is a reporting
convenience --- the command-line summary prints such a dimension as
\code{LOW} (\code{runner.py:157-163}) --- not an exclusion.

\code{hard\_floor} is declared by eleven of the composite dimensions and is
never read anywhere in the package. It is a dormant declaration; see
\S\ref{sec:absences}.

\subsection{The verdict}

The verdict is a label attached to the composite, computed by comparing the
composite against three thresholds in order (\code{runner.py:24-31}):

\begin{longtable}{@{}p{4.4cm}p{4.4cm}@{}}
\toprule
Composite at or above & Label \\
\midrule
\endhead
$0.70$ & \code{VIABLE} \\
$0.50$ & \code{BORDERLINE} \\
$0.30$ & \code{REGRESSED} \\
(below all of the above) & \code{COLLAPSED} \\
\end{longtable}

The verdict is \emph{not} a gate: it fails nothing and is counted nowhere. It
exists to be read by a human.

\section{Gates: binary invariants, reported apart}

\subsection{What makes something a gate}

A gate is a scorer whose kind is \code{"gate"}
(\code{mal-core/src/mal\_core/scoring/scorers/base.py:56-64}). Gates check
data hygiene and input validity, not biology. Three properties distinguish
them from composite dimensions, and each is visible in the code:

\begin{itemize}
\item \textbf{They never enter the composite.} The runner routes results by
      kind (\code{runner.py:103-106}), and the gate classes declare a weight
      of zero (\code{gates.py:44}, \code{:85}, \code{:121}, \code{:152},
      \code{:202}, \code{:257}).
\item \textbf{A failure is a warning, not a penalty.} The gate module's
      docstring says it: a failed gate is a visible warning in the scorecard
      but never changes the composite (\code{gates.py:1-7}).
\item \textbf{Missing data is a skip, not a failure.} Every gate returns a
      skipped result when its artefacts are absent, on the stated principle
      that missing data is not a data-integrity failure (\code{gates.py:50-51},
      \code{:91-95}, \code{:127-131}, \code{:158-169}, \code{:213-216},
      \code{:266-273}). A skipped gate appears in \code{skipped}, not in
      \code{gates}.
\end{itemize}

\subsection{The six gates}

\begin{longtable}{@{}p{2.4cm}p{5.2cm}p{6.4cm}@{}}
\toprule
Gate & What it checks & Rule and units \\
\midrule
\endhead
\code{G6\_data\_valid} & Every band value of every state COG is finite and
non-negative & Fraction of valid values over all bands of all state files;
the gate passes only at a perfect fraction. Above 48 state files a uniform
sample of 48 is taken, the full sweep being a test concern
(\code{gates.py:52-54}). \\
\code{G7\_determinism} & Per-seed state files are byte-identical & Files
grouped by the seed token in the filename; every group with two or more
members is compared by MD5; the score is one when no group mismatches
(\code{gates.py:89-113}). \\
\code{G10\_perf} & Wall time of the run & Read from \code{wall\_s} in
\code{run\_metadata.json}; the gate passes at or below 30 s, and the score
decays as $30/w$ otherwise (\code{gates.py:125-144}). \\
\code{G12\_host\_input} & The input human grid total matches a census
reference & Total of the \code{human} variable of the AOI's host raster,
excluding the NoData sentinel \mskey{ingest.nodata_sentinel}, divided by a
reference total that defaults to $27\,000\,000$ and can be overridden
through the context; the gate passes when the ratio lies in $[0.8, 1.2]$
(\code{gates.py:156-194}). \\
\code{G14\_mobility} & Mobility origin--destination matrices are
row-stochastic & Each \code{.csr} file next to the run directory (or in its
parent, or in a sibling \code{hosts} directory) is read as a sparse matrix;
the fraction of rows whose values sum to one within a tolerance of $0.01$ is
computed per file, the minimum over files is the score, and the gate passes
at or above $0.99$ (\code{gates.py:206-245}). \\
\code{G24\_urban\_ratio} & Urban cells carry less adult density than rural
ones & Mean adult density over urban-class $30$ cells divided by the mean
over rural-class $50$ cells of the final state; the gate passes at or below
$0.7$ (\code{gates.py:261-303}). The gate is \emph{manual}: it needs the
AOI's urban-class raster. \\
\end{longtable}

\subsection{Why the data-validity rule is finite and non-negative}

Under contract \mskey{contracts.cogs_state} the two state bands are raw
per-cell adult counts --- \code{adult\_patch\_count} and
\code{adult\_cell\_count} (\code{mal-core/src/mal\_core/abm/src/main.cpp:386-390},
and the contract table above) --- so a value above one is an ordinary adult
count and not by itself an error. The rule the gate enforces is therefore
\emph{finite and $\geq 0$} (\code{gates.py:35-40}): it rejects non-finite
values, the NoData sentinel and negative values, and passes only when
every value of every band of every sampled state file is valid
(\code{gates.py:52-54}).

What the code does \emph{not} record is the history of that rule. The
gate's docstring asserts that the contract changed and that the earlier
$[0,1]$ envelope no longer applies (\code{gates.py:35-40}); the tree shows
only the current contract version, so that assertion cannot be checked
against the code and is reported here as the module's own account of its
motivation.

This has a consequence for the dimensions that read the same bands, stated
in \S\ref{sec:dimensions}.

\section{The thirteen composite dimensions}
\label{sec:dimensions}

\subsection{What each one measures}

\begin{longtable}{@{}p{2.5cm}p{3.4cm}p{1.9cm}p{3.3cm}p{2.9cm}@{}}
\toprule
Dimension & Quantity & Artefact & Units and target band & Tier (this document's
classification) \\
\midrule
\endhead
\code{D1\_habitat\_occupancy} & Fraction of habitat cells holding at least
one adult in the final snapshot & final state COG, bands 1 and 2 &
dimensionless; band $0.20$--$0.70$; Gaussian falloff around $0.45$ of width
$0.15$ & literature-anchored (Bogh 2004; Ndenga 2011) \\
\code{D2\_survival} & Mean adult lifespan implied by the steady-state daily
adult survival rate & cohort log & days; window $8$--$15$ d; Gaussian around
$11$ d of width $4$ d, with no plateau at the ends of the window &
literature-anchored (Saarman 2019; Midega 2007) \\
\code{D3\_eip} & Fraction of infected mosquitoes whose extrinsic incubation
period has completed, at day 30 & cohort log & dimensionless; band
$0.20$--$0.50$; Gaussian around $0.35$ of width $0.15$ &
literature-anchored (Mordecai 2013) \\
\code{D4\_stability} & Adult population trajectory: collapse ratio and
explosion ratio & cohort log & ratio to the day-0 adult count; collapse band
$0.10$--$0.80$, explosion at or below $5\times$; the score is the mean of
the two sub-scores & engineering prior \\
\code{D5\_morans} & Moran's $I$ of the occupied cells of the final snapshot,
queen contiguity, on active pixels only & final state COG, band 1 &
dimensionless; band $0.20$--$0.85$; Gaussian around $0.5$ of width $0.25$ &
literature-anchored (spatial clustering of larval habitats) \\
\code{D8\_coupling} & Pearson correlation between the two state bands over
the cells where band 1 is positive & final state COG, bands 1 and 2 &
correlation coefficient; target $r \geq 0.30$; the score is $r/0.50$ capped
at one & engineering prior \\
\code{D11\_aquatic\_structure} & Mean larval share and mean pupal share of
the aquatic population over the last 30 aquatic files & per-day aquatic
JSONs & dimensionless; larva band $0.65$--$0.90$, pupa band $0.02$--$0.15$;
the score is the mean of the two sub-scores & literature-anchored (Bayoh \&
Lindsay 2003; Ouedraogo 2024) \\
\code{D15\_persistence} & Long-horizon persistence: survival against the
population peak, absence of extinction, spatial expansion & per-day aquatic
JSONs plus state COGs & dimensionless; three sub-scores combined by a
geometric mean & engineering prior \\
\code{D16\_detection\_coverage} & Fraction of the AOI's real detection spots
that the simulated population reached & final state COG plus the AOI's
occurrence table & dimensionless; target at or above $0.50$ at a radius of
$5$ km, with sensitivity radii of $2$, $5$ and $10$ km reported in the
diagnostics & scenario default (launch-and-measure) \\
\code{D19\_pool\_persistence} & Whether the proxy-habitat model keeps
pluvial availability alive after a rain pulse and keeps a positive dry-season
floor & \emph{nothing from the run} (\S\ref{sec:contractdims}) &
dimensionless; four checks worth $0.3$, $0.2$, $0.2$ and $0.3$ & model
definition \\
\code{D20\_washout\_response} & Whether heavy rain produces a washout of
aquatic cohorts with the expected shape & \emph{nothing from the run}
(\S\ref{sec:contractdims}) & dimensionless flush fraction; four checks worth
$0.3$, $0.3$, $0.2$ and $0.2$ & model definition \\
\code{D21\_dispersal\_kernel} & Median and 90th percentile of the
displacement of the adult population from the seeding centroid & state COGs
at the first snapshot and at the latest snapshot at or before day 30 &
kilometres; median band $0.15$--$0.70$ km, 90th-percentile band
$0.70$--$3.00$ km; the score is the mean of the two sub-scores &
literature-anchored (Costantini 1996; Thomas 2013; Midega 2007) \\
\code{D25\_cases\_ghana} & Correlation between the simulated monthly
incidence climatology and the observed reported-cases climatology &
transmission daily log plus the pinned DHIMS-II table & correlation
coefficient; target $r \geq 0.60$ & literature-anchored (Seidu 2026) \\
\end{longtable}

\subsection{How a measurement becomes a score}

The table above states the band each dimension compares against. The
mechanisms differ, and the difference is the difference between the six
shapes listed in \S\ref{sec:composite}: a plateau-and-falloff (D1, D3, D5,
D25), the same shape averaged over two bands (D4, D11, D21), a Gaussian
without a plateau (D2), a ratio against a target (D8, D16), a geometric
mean of three sub-scores (D15), and a weighted sum of boolean checks
(D19, D20).

Two behaviours change what a score means and are worth stating separately:

\begin{itemize}
\item \textbf{D3 skips rather than punishes a degenerate run.} If the
      completed-EIP fraction is constant across the cohort log --- the state
      of a run made without the transmission module --- the dimension
      reports a skip (\code{d3\_eip.py:44-50}).
\item \textbf{D16 distinguishes extinction from zero coverage.} A run whose
      final state has no occupied cell scores zero, carrying an
      \emph{extinct} diagnostic when the cohort log shows zero adults at the
      last day (\code{d16\_detection\_coverage.py:102-117},
      \code{:145-157}).
\end{itemize}

\subsection{Where the state bands are read, and the mismatch that is there}

Three dimensions read the two bands of the state COG, and their docstrings
describe them in the vocabulary of an \emph{earlier} contract: band 1 as
adult density, band 2 as aquatic-capable habitat or host-seeking pressure
(\code{d1\_habitat\_occupancy.py:3-5}, \code{d5\_morans.py:1-4},
\code{d8\_coupling.py:1-6}). Under contract 4.0 both bands are raw
per-cell adult counts: band 1 counts the adults associated with their patch
after dispersal, band 2 counts the adults at their current position
(\code{mal-core/src/mal\_core/abm/src/main.cpp:386-390}).

The consequence is stated plainly rather than smoothed over: \emph{as
implemented today, these three dimensions threshold, count or correlate
adult counts, not habitat suitability}.

\begin{itemize}
\item \textbf{D1} divides the number of cells where band 1 is positive by
      the number of cells where band 2 is positive
      (\code{d1\_habitat\_occupancy.py:51-61}). With both bands being adult
      counts, this is the share of adult-holding cells that also hold adults
      under a second, differently-defined count; the scorer's own skip
      condition fires when band 2 has no positive cell at all
      (\code{d1\_habitat\_occupancy.py:58-59}).
\item \textbf{D5} computes Moran's $I$ on band 1
      (\code{d5\_morans.py:70-71}), which is the spatial autocorrelation of
      adult presence --- a meaningful quantity, but not the habitat
      clustering quantity the docstring names.
\item \textbf{D8} correlates band 1 against band 2 on the cells where band 1
      is positive (\code{d8\_coupling.py:32-47}), that is, it measures how
      strongly the two adult-count definitions agree. Because both are
      counts of the same population under two spatial definitions, a strong
      positive correlation is expected by construction, and the dimension
      cannot fail for the reason it was written to test.
\end{itemize}

This is a defect in the measurement layer, not in the contract. It is
recorded here as an open item (\S\ref{sec:absences}) rather than papered
over. A consumer reading a D1, D5 or D8 score from a contract-4.0 run should
read it as a statement about adult counts, and should not read it as a
statement about habitat suitability.

\subsection{Two dimensions that score the model contract instead of the run}
\label{sec:contractdims}

D19 and D20 accept a run directory and a context and then ignore both
(\code{d19\_pool\_persistence.py:51-71},
\code{d20\_washout\_response.py:41-63}). They exercise the habitat
contract's own rules and check them against thresholds. The two differ in
how they do it: D19 writes the pluvial availability recursion out inline,
as a scalar step on a single cell (\code{d19\_pool\_persistence.py:65-71}),
whereas D20 imports and calls the shared washout-fraction function
(\code{d20\_washout\_response.py:20}, \code{:46}, \code{:51}, \code{:56},
\code{:61}) --- so D20 exercises the shipped implementation, while D19
exercises a transcription of the recursion written inside the scorer. The
constants they
use are the shared model constants: the pluvial rain scale is
\mskey{pluvial.rain_scale_mm} and the drying memory is
\mskey{pluvial.dry_memory_days}, the washout threshold is
\mskey{aquatic_stress.wash_threshold_mm} and the washout cap is
\mskey{aquatic_stress.wash_fraction_max}
(\code{d19\_pool\_persistence.py:28-29},
\code{d20\_washout\_response.py:24-25}).

Their scenarios are fixed. D19 applies a $30$ mm pulse on three consecutive
days and then ten dry days, with an activation floor of $0.05$
(\code{d19\_pool\_persistence.py:78-90}); D20 evaluates the washout rule at
$30$, $40$, $50$ and $80$ mm (\code{d20\_washout\_response.py:45-63}).

Their provenance is therefore \emph{model definition}: they assert that the
contract's own rules still behave as the contract says. Their practical
effect on the composite is that two of the thirteen dimensions are
insensitive to the run being scored --- they can only contribute their full
weight or collapse it, and they contribute the same value to every run.

\section{The AOI convention: explicit in, resolved by convention}

\subsection{The slug is never inferred}

The AOI slug reaches a scorer only through an explicit \code{--aoi} on the
command line (\code{mal-core/src/mal\_core/cli.py:413}), which the runner
copies into the context (\code{runner.py:34-40}). A scorer that needs the
slug asks for it through \code{require\_aoi()}, which returns \code{None}
rather than guessing (\code{scorers/base.py:164-172}). The reason is in the
code: run outputs may live in any user-chosen directory, so guessing the
AOI from paths or from run metadata is forbidden by convention.

This diverges from the command-line help text, which says the AOI is
auto-detected from run metadata if omitted (\code{cli.py:413}). The code is
authoritative; the help text is wrong, and it is listed in
\S\ref{sec:absences}.

\subsection{Where per-AOI inputs come from}

The \emph{slug} is explicit; the \emph{path} to that AOI's data is resolved
by convention --- repository-relative glob patterns with
\texttt{\{aoi\}} substituted, first existing hit wins
(\code{scorers/base.py:204-223}). The patterns in use are those in the
artefact table above: the occurrence table for D16, the host grid for G12
and the urban-class raster for G24. D25 is the exception: its dataset is a
fixed, pinned public file rather than something the AOI provides. The
scorer downloads it on demand into \code{data/ghana/scoring/} and refuses
to use a copy whose SHA-256 checksum does not match the pinned value
(\code{d25\_cases\_ghana.py:39-77}); it also refuses any AOI other than
Ghana, because the table covers three of Ghana's regions
(\code{d25\_cases\_ghana.py:192-197}).

\subsection{The grid transform, and why it is not trivial}

Two dimensions (D16, D21) need to convert between grid cells and
longitude/latitude. The helper \code{grid\_transform()} trusts the
GeoTIFF's own affine when the grid is north-up (the off-diagonal terms
vanish), and otherwise falls back to the \code{transform} field of the run
summary JSON (\code{scorers/base.py:175-201}). The helper's own docstring
gives its reason: the engine's state COGs have historically stored a
garbled affine, with GDAL ordering fed to an affine constructor
(\code{scorers/base.py:175-182}). That historical claim is the module's
account of itself --- no code in the tree records the past behaviour, and
the current writer does not show it --- so it is reported here as an
assertion, not as a verified fact. When neither source yields a
usable transform, both dimensions skip rather than guess
(\code{d16\_detection\_coverage.py:79-81},
\code{d21\_dispersal\_kernel.py:71-73}).

\section{The optional judge: a qualitative reading, not evidence of the same
kind}

\subsection{What it does}

The judge posts the finished scorecard JSON to a chat-completions endpoint
on OpenRouter together with a bundled system prompt, and validates the
reply against a schema (\code{mal-core/src/mal\_core/scoring/llm\_scorer.py:74-95}).
The schema requires a verdict from a four-value vocabulary, a composite
estimate in $[0,1]$, and lists of concerns, recommendations and literature
grounding (\code{llm\_scorer.py:82-103}).

\subsection{Why it is not the same kind of evidence as a deterministic scorer}

\begin{itemize}
\item \textbf{It cannot gate anything.} The runner calls it after the
      composite and the verdict are already computed, and attaches its reply
      only under the optional key \code{verdict\_llm}; if the call fails the
      key is omitted entirely (\code{runner.py:135-142}). The deterministic
      composite and verdict are unaffected in every path.
\item \textbf{It can be absent.} Three failure paths --- no API key, a
      schema-validation failure, any transport or other exception --- all
      produce a fallback dictionary carrying \code{verdict="unknown"} and a
      human-readable reason, and no exception escapes to the caller
      (\code{llm\_scorer.py:222-254}, \code{:296-308}, \code{:359-368}). The
      scorecard is written either way.
\item \textbf{It is reproducible only in the caching sense.} It is run at
      temperature zero and the request is content-hashed and cached on disk,
      so repeated runs over an unchanged scorecard are free and return the
      same answer (\code{llm\_scorer.py:111-131}, \code{:296-303},
      \code{:319-321}). The answer is still the output of a third-party
      model on a prompt, not a function of the run.
\item \textbf{It is judging a stale description of the scorecard.} The
      bundled prompt still describes a ten-dimension scorecard and reads
      band 1 as larval habitat saturated at one and band 2 as a normalised
      adult density (\code{scoring/prompts/llm\_verdict.md:3-4},
      \code{:71-79}), whereas the registry now holds thirteen composite
      dimensions and the state bands are raw adult counts. The judge's
      context is therefore partly wrong.
\end{itemize}

\subsection{How it is kept cheap}

The request payload --- the scorecard JSON plus the model identifier --- is
hashed with SHA-256, and the verdict is cached under a content-addressed
file in a cache directory (\code{llm\_scorer.py:111-131},
\code{:177-187}). The cache is best-effort: a failed write is swallowed
(\code{llm\_scorer.py:208-222}). The API key is read from the environment or
from a \code{.env} file walked up from the package directory
(\code{llm\_scorer.py:134-174}). The default model identifier, the base URL
and the per-request timeout are module constants
(\code{llm\_scorer.py:44-52}).

\section{Measurement campaigns: the two commands that build a run first}
\label{sec:campaigns}

\subsection{Detection coverage}

The \code{validate-detections} command seeds mosquitoes at the AOI's real
detection points, runs the model for a year by default, and then scores
with the detection-coverage dimension explicitly enabled
(\code{mal-core/src/mal\_core/cli.py:466-556}, \code{:551-555}). The
occurrence table is read with the shared loader, which deduplicates the
detection spots (\code{scorers/base.py:233-254}). The dimension then asks
how many of those spots the simulated population actually reached: a spot
counts as covered when at least one occupied cell lies within the radius of
the spot in the final snapshot (\code{d16\_detection\_coverage.py:87-126}).

\subsection{Reported cases}

The \code{validate-cases} command runs a multi-year simulation with a
uniform seeding across candidate patches and the whole vector pool
initially infected, discards a warm-up, and then scores with the
reported-cases dimension explicitly enabled
(\code{cli.py:559-639}, \code{:633-638}). The dimension compares the
simulated calendar-month climatology against the observed one; when the
simulation's last twelve blocks fall in a single calendar year that the
dataset also covers, the comparison is made year by year, otherwise it falls
back to the pooled climatology (\code{d25\_cases\_ghana.py:216-246}).

Both campaigns exist because their dimensions are only meaningful for a run
constructed in that particular way; scoring an arbitrary run with them
would measure the wrong thing.

\subsection{The feedback report}

The feedback step renders a Markdown report from a scorecard: the composite
and its verdict, the failed gates, the skipped scorers, the three weakest
dimensions, and --- when a baseline scorecard is supplied --- the composite
delta and the dimensions that regressed by more than $0.05$
(\code{mal-core/src/mal\_core/scoring/feedback.py:16-53}). The command line
re-scores the run directory and, optionally, a baseline directory, and
prints the report (\code{cli.py:841-861}).

\subsection{A worked situation}

An operator has changed a parameter in the model, produced a new run
directory, and wants to know whether the change helped. The operator asks
for a scorecard of the new run, and of an older baseline run, and then asks
for feedback with the baseline attached. What happens is: the new run is
scored dimension by dimension; the composite is a weighted geometric mean
over the dimensions that could be measured; the verdict label follows from
the composite; any dimension whose artefacts were absent appears in the
skipped list with a reason instead of dragging the composite down; the
feedback report then lists the three weakest dimensions and names any
dimension whose score fell by more than $0.05$ relative to the baseline. The
operator sees whether the change moved the weakest part of the model rather
than merely moving the average.

\section{Declarations of absences}
\label{sec:absences}

\subsection{What is not implemented}

\begin{itemize}
\item \textbf{There is no pytest calibration harness in this tree.} No
      \code{calibration/} directory, no threshold table, no
      diff-versus-best tracker: the stage's base module states that the old
      harness and its duplicated scorers were deleted, and the thresholds it
      used to hold now live inside the scorer modules or, where they belong
      to the model, in the shared configuration
      (\code{mal-core/src/mal\_core/scoring/scorers/base.py:1-6}).
\item \textbf{The composite does not gate anything.} No code in the tree
      consumes the composite or the verdict as a promotion criterion; the
      \code{score} command computes, writes and prints, and stops
      (\code{mal-core/src/mal\_core/cli.py:409-463}). Its only non-zero exit
      is the missing-\code{--run-dir} argument error
      (\code{mal-core/src/mal\_core/cli.py:448-450}), which is unrelated to
      the verdict. Whether an external process is meant to use the verdict
      is not settled by the code.
\item \textbf{The judge has no mechanical effect.} Its output is an optional
      key (\code{runner.py:135-142}).
\end{itemize}

\subsection{What is declared but dormant}

Three class attributes are declared and read by no code in the package:

\begin{longtable}{@{}p{3.2cm}p{4.4cm}p{6.4cm}@{}}
\toprule
Attribute & Declared by & Status \\
\midrule
\endhead
\code{weight} & every scorer (\code{base.py:68}) & The runner asks the
configuration for the weight, never the class (\code{runner.py:112},
\code{config.py:30-31}), and the listing also reads the registry table
(\code{registry.py:116}). The class values agree with the table today, so
nothing is wrong now, but there are two sources of truth and no rule about
which wins when they disagree. \\
\code{hard\_floor} & eleven composite dimensions & Never read anywhere in
the package. Its intended role is not settled by the code. \\
\code{aoi\_generalizable} & five of the six gates (\code{gates.py:46},
\code{:87}, \code{:123}, \code{:154}, \code{:204}) & Never read
anywhere in the package. The sixth gate, the urban-ratio gate, declares no
such attribute at all; it carries \code{manual = True}
(\code{gates.py:259}) instead. \\
\end{longtable}

Two further items deserve a mention. The Mann--Whitney area helper
\code{auc\_rank()} is defined in the base module (\code{base.py:257-285})
and called by no dimension --- a leftover of the retired occurrence-AUC
scorers. The observation loader \code{load\_occurrence\_spots()}
(\code{base.py:233-254}) is \emph{not} a leftover: it is used by the
detection-coverage dimension (\code{d16\_detection\_coverage.py:83}) and by
the detection campaign command (\code{cli.py:497-512}), so it is part of
the stage's contract.

Two dimensions (D19, D20) never inspect the run they are asked to score;
they score the habitat contract in isolation and contribute to the
composite as if they were measurements of the run
(\S\ref{sec:contractdims}). Their thresholds are also written literally
instead of being taken from the declared minimum score: both compare their
score against a literal $0.70$ (\code{d19\_pool\_persistence.py:144},
\code{d20\_washout\_response.py:83}), which coincides with their declared
minimum score but breaks the uniform pattern.

The judge is prompted against a stale description of the scorecard: the
bundled prompt describes ten dimensions and pre-contract band semantics
(\code{scoring/prompts/llm\_verdict.md:3-4}, \code{:71-79}). The module's
docstrings, by contrast, agree with its implementation: they describe a
direct HTTP call with no agent harness and no agent loop
(\code{llm\_scorer.py:16-19}), and the entry point's docstring says it
posts a JSON payload to the OpenRouter \code{chat/completions} endpoint
with \code{requests} and parses the reply into a \code{Verdict}
(\code{llm\_scorer.py:269-277}) --- which is what the code does
(\code{llm\_scorer.py:330-353}).

Finally, the command-line help text for \code{--aoi} promises an automatic
detection that the code forbids (\code{cli.py:413} against
\code{scorers/base.py:164-172}).

\section{What is not answered by the code}

The following questions were posed for this stage and the implementation does
not settle them. They are listed so that a reader does not mistake silence
for a decision.

\begin{enumerate}
\item \textbf{What do the state bands mean, as the scorers read them?} The
      engine writes two raw adult-count bands, and three dimensions read
      them with the vocabulary of a retired contract
      (\S\ref{sec:dimensions}). Whether those three dimensions should be
      re-specified or retired is not decided by the code.
\item \textbf{What is \code{hard\_floor} for?} It is declared and never read.
      Whether it is meant to be a hard rejection threshold, a reporting band,
      or residue of the deleted harness is not decided by the code.
\item \textbf{Which weight source is canonical?} The registry table and the
      class attribute agree today; the code does not say what happens when
      they do not.
\item \textbf{Is the composite meant to be a promotion criterion?} Nothing in
      the tree enforces it.
\item \textbf{Is the judge's prompt an in-flight fix or a settled
      description?} It disagrees with the current scorecard.
\item \textbf{Should the contract-testing dimensions belong in the
      composite?} Two of the thirteen dimensions score the model contract
      rather than the run, and their presence narrows what the composite
      means.
\end{enumerate}

\section{Evidence appendix}

\subsection{Implementation notes a replicator would need}

\begin{itemize}
\item Artefacts are found by filename pattern, not by a manifest. A run that
      names its snapshots differently is scored as an empty run, and every
      dimension skips.
\item The state-COG fallback chain is ordered, and the first pattern that
      matches wins: per-day snapshots, then per-seed snapshots, then a
      single final raster (\code{scorers/base.py:89-102}).
\item The run summary is the first JSON file in the directory whose content
      carries a contract version or an AOI slug, excluding files whose names
      carry \emph{cohort}, \emph{aquatic}, \emph{day} or \emph{scorecard}
      (\code{scorers/base.py:142-161}).
\item A dimension that raises is recorded as skipped with the exception text,
      and the composite is unaffected (\code{runner.py:96-98}).
\item Scores are rounded to four decimals by each scorer before they enter
      the composite (each scorer's return statement); the composite itself
      is rounded to four decimals when it is written into the scorecard
      (\code{runner.py:123}).
\item The NoData sentinel is not a literal in the scoring code: it is read
      from the shared habitat configuration (\code{scorers/base.py:21-22}),
      where it is \mskey{ingest.nodata_sentinel}.
\end{itemize}

\subsection{Provenance of the numbers in this stage}

\begin{itemize}
\item \textbf{Model definition.} The pluvial rain scale and drying memory,
      the washout threshold and cap
      (\code{d19\_pool\_persistence.py:28-29},
      \code{d20\_washout\_response.py:24-25}), and the NoData sentinel are
      read from the shared model configuration through its loader.
\item \textbf{Literature-anchored.} The occupancy bands (Bogh 2004; Ndenga
      2011), the lifespan window (Saarman 2019; Midega 2007), the incubation
      window (Mordecai 2013), the aquatic stage ratios (Bayoh \& Lindsay
      2003; Ouedraogo 2024), the dispersal percentiles (Costantini 1996;
      Thomas 2013; Midega 2007) and the reported-cases target (Seidu 2026)
      are literals inside the scorer modules, and each module's docstring
      names the source. The registry module's docstring lists anchors for
      eight of the thirteen dimensions only --- D1, D2, D3, D11, D16, D19,
      D21 and D25 (\code{registry.py:7-18}) --- and of those, only D21's
      anchor names a file under the repository's \code{papers/} directory.
      D4, D5, D8, D15 and D20 have no anchor in that docstring; the anchors
      they do carry are in their own modules (D5 cites Tokarz \& Novak 2018
      and Bissett 2026, D19 cites Bøgh 2004 and Ndenga 2011, which differ
      from the registry's own anchor for D19 --- Depinay 2004 and Ward 2015
      --- \code{registry.py:13-14}). The registry's opening line, that all
      papers live in \code{papers/}, is a statement about the repository,
      not a citation list.
\item \textbf{Engineering prior.} Numbers with no literature behind them,
      written straight into the scorer: the activation floor and pulse of
      D19, the sampling cap of G6, the wall-time target of G10, the
      reference population of G12, the tolerance and pass fraction of G14,
      the urban-ratio limit of G24, the coverage radii of D16.
\item \textbf{Scenario default.} The choices the two campaigns make: the
      detection campaign's run length, the number of simulated years and
      the warm-up of the reported-cases campaign (\code{cli.py:466-556},
      \code{cli.py:559-639}).

The assignment of a given literal to the engineering-prior tier or to the
scenario-default tier is this document's classification and not something
the code makes: both kinds of value are plain literals in the module that
uses them, with no marker distinguishing them
(\code{d16\_detection\_coverage.py:43-45},
\code{d4\_stability.py:56-66}).
\end{itemize}

\subsection{What this document's numbers are}

Every value in the generated tables --- the composite weights, the minimum
scores, the hard floors, the gate descriptions and the output contracts ---
comes from \code{docs/concepts/generated/params-scoring.tex} and
\code{docs/concepts/generated/contracts.tex}, which are emitted at build
time from the model configuration and the contract registry.

Every other number in this document is a literal inside a scorer module,
cited to the file and line where it is written. These are \emph{readings of
the tree}: they are verified against the source as quoted, they are not
generated, and they will silently go stale if the scorer literals change
without this document being updated. That is the price of literals living
outside the shared configuration, and it is stated here so that a reader
knows which numbers to re-check when a scorer changes.

\subsection{What is verified and what is not}

\begin{itemize}
\item \textbf{Verified by reading the source.} Every behavioural claim above
      resolves to a file and a line a reviewer can open: the order of
      operations, the composite formula and its degenerate cases, the gate
      semantics, the treatment of failures and skips, the judge's
      non-gating status, the artefact discovery patterns, the AOI convention
      and the dormant declarations.
\item \textbf{Not verified by execution.} This document reports the output of
      no real scoring run. No composite value, no dimension score and no
      gate outcome is asserted here, because no delivered run artefact set
      was measured for this document; any number that would have required
      such a measurement is therefore absent.
\item \textbf{Known to be stale in the implementation.} The judge's prompt,
      the command-line help text for the AOI option, and the band semantics
      that three dimensions still assume.
\end{itemize}

\end{document}
